%% Conclusion: What the Program Has Shown

This thesis began with a question that felt philosophical: what must
an observer be, for coherent probabilistic experience to be possible?
It ends with an algorithm.  The distance between those two points is
the measure of what the program has established.

\section*{The arc}

Chapter~\ref{chap:paper-minus-one} showed that an observer whose
discriminative acts satisfy finite Boolean closure, and who is
committed to acting coherently toward a world that exceeds their
current distinctions, is forced to possess a $\sigma$-algebra.  Not
as an assumption.  As the form that coherent commitment must take.

Chapter~\ref{chap:paper-zero} showed that when a compatible family
of such observers --- queries organized by a refinement order ---
satisfies collective exhaustion and a realizability condition on the
projective limit, their local commitments assemble into a unique
global probability measure $P$ on the realization space $\Omega$.
Nothing is hidden: $P$ is entirely determined by the observable
interface.

Chapter~\ref{chap:paper-one} showed that the predictive content of
each query factors through a minimal sufficient statistic: the
minimal predictive query $Q_*$.  Prediction is structured, not
arbitrary.

Chapter~\ref{chap:paper-two} showed that as time advances, the
predictive structure evolves as a semigroup.  The dynamics are
coherent across time in the same way that the queries are coherent
across refinement levels.

Chapter~\ref{chap:paper-three} showed that this semigroup can be
instantiated via delay queries and approximated from finite data.
The abstract program becomes computable.

\section*{What was derived, what was assumed}

The program is honest about the distinction.  Three things were
derived: the $\sigma$-algebra (from coherent commitment, under the
witnessing conjecture of Paper~$-1$), $\sigma$-additivity (from
tightness or perfectness, for the architectural cases), and the
predictive semigroup (from observable compatibility alone).

Three things remain as structural hypotheses: upper-directedness
(SP4: assumed, with an operational motivation but no derivation),
realizability of $\Omega$ (SP3: assumed, independence from
compatibility unknown), and the general case of SP2 beyond compact
and standard Borel outcome spaces.

These are honest gaps.  The program does not claim to derive
everything from nothing.  It claims to derive more than was
previously known to be derivable, and to identify precisely what
the remaining assumptions are and what they mean.

\section*{The open frontier}

The central open question is the witnessing conjecture of
Paper~$-1$: that sequential upper-directedness provides, for every
observer at every finite level, a witness from above that forces
continuity at $\emptyset$.  A proof would close SP1 and give SP2
a topology-free route simultaneously, subsuming the compact and
standard Borel cases as special instances of a single algebraic
result.

If that conjecture is true, the program's chain reads cleanly:
\[
\text{bounded discriminability}
\;\to\;
\text{refinement}
\;\to\;
\text{coherence}
\;\to\;
\text{probability}
\;\to\;
\text{prediction}
\;\to\;
\text{dynamics}
\;\to\;
\text{computation.}
\]
Every arrow forced.  None assumed.  That is the program's claim.
It is not yet fully proved.  But the ground has been laid, the
gaps are named, and the next steps are clear.
